\section{One simultaneous window for the two-primary and reference-layer terms}
\label{sec:simultaneous-cohomological-windows}

\status{TW1--TW8 reconstruct the complete supplier
$\exists K\,\forall C_t\ge1,\ \texttt{ThmBab}\ K\ C_t$.
The input from the quantitative cohomology branch is the literal
$\exists C_K,\ \texttt{LemmaK}\ C_K$, supplied by KB8 with $C_K=31$;
finite Poitou--Tate duality
is supplied by LD3 and its actual local-field hypotheses by LK2--LK3.
Ferrero--Washington is used in its classical $\mu=0$ form, and its
translation to the precise uniform class-rank statement in the source
is proved below.  Every tower, exceptional-twist, common-window and
field-model step is included.  This chapter does not assume
\texttt{FixedIn}, \texttt{TwoIn}, \texttt{ABound}, \texttt{FirstBound}
or \texttt{ThmBab} as an additional hypothesis.}

For a number field $F$, a prime $\ell$ and $k\ge1$, put
\[
 A_\ell(k)=(\mathbb Q_\ell/\mathbb Z_\ell)(-k),\qquad
 s_\ell(F,k)=\operatorname{length}_{\mathbb Z_\ell}
 \ker\!\left(H^1(G_{F,S_\ell},A_\ell(k))
       \longrightarrow\prod_{v\mid\ell}H^1(F_v,A_\ell(k))\right).
 \tag{TW.1}
\]
Here $S_\ell$ consists of exactly the primes above $\ell$.
All cochains with these torsion coefficients are continuous for the
discrete topology.  Thus $s_\ell(F,k)$ is the source's
\texttt{shaInv F l k}, and also \texttt{shaLen D k} for every
algebraic closure and every admissible local-data record $D$ on the
given $F$.  At first, lengths may be infinite.  The compact dual
coefficient is $T_\ell(k)=\mathbb Z_\ell(k+1)$, and
\[
 V_\ell(F,k;D)=
 \operatorname{length}_{\mathbb Z_\ell}H^2(G_{F,S_\ell},T_\ell(k))
 \tag{TW.2}
\]
is \texttt{VLen D k}.

Write $\operatorname{VV}_\ell(F,t,B)$ for the assertion, for every
such closure and $D$, that all the lengths in (TW.1) with
$1\le k<t$ are finite and that their sum is at most $B$.
This is exactly \texttt{VVat}; at $\ell=2$ it is
\texttt{BoundedHeightTwoPlaceCondition}.
An \emph{$m$th cyclotomic layer} of $F$ is a field
$F(\mu_{\ell^a})$ of degree $\ell^m$ over $F$ for some $a$,
allowing any isomorphic field model equipped with its $F$-algebra
structure.  This is the literal \texttt{IsLayer l m F F'}.

\subsection{Named classical inputs}

Only the following classical inputs are retained in this chapter;
their finite-layer, coefficient and quantitative consequences are
proved in TW1--TW8.

\begin{description}[style=nextline,leftmargin=0pt,labelwidth=0pt,labelsep=0pt]
\item[TW-FW: Ferrero--Washington.]
For every finite abelian number field $F/\mathbb Q$ and every prime
$\ell$, the unramified Iwasawa module of its cyclotomic
$\mathbb Z_\ell$-extension has $\mu=0$.
This includes $\ell=2$.
See B.~Ferrero and L.~C.~Washington,
\emph{The Iwasawa invariant $\mu_p$ vanishes for abelian number fields},
Annals of Mathematics \textbf{109} (1979), 377--395,
\url{https://doi.org/10.2307/1971116},
and the precise all-prime formulation in R.~Sharifi,
\emph{Iwasawa Theory}, Theorem 6.2.1,
\url{https://math.ucla.edu/~sharifi/notes/iwasawa-ch06.html}.

\item[TW-IW: the ordinary Iwasawa module theorem.]
For a $\mathbb Z_\ell$-extension $F_\infty/F$, the inverse limit
$X=\varprojlim\operatorname{Cl}(F_n)\{\ell\}$ under norms is a
finitely generated torsion
$\Lambda=\mathbb Z_\ell[[\operatorname{Gal}(F_\infty/F)]]$-module.
Artin reciprocity identifies it with the Galois group of the maximal
unramified abelian pro-$\ell$ extension of $F_\infty$.
See Sharifi, Corollary 3.2.7 and Proposition 3.2.11,
\url{https://math.ucla.edu/~sharifi/notes/iwasawa-ch03.html}.
The structure theorem used with $\mu=0$ is Theorem 2.2.13,
together with Remark 2.2.11 and Lemma 2.2.12,
\url{https://math.ucla.edu/~sharifi/notes/iwasawa-ch02.html}.
It says that a torsion $\Lambda$-module has, up to finite kernel and
cokernel, distinguished-polynomial and $\ell$-power factors;
$\mu=0$ removes the latter.

\item[TW-CF: finite class-field reciprocity.]
The maximal everywhere unramified abelian extension of a totally
complex number field has Galois group its ideal class group.
Artin maps intertwine restriction of automorphisms with ideal norms.
See J.~S.~Milne, \emph{Class Field Theory}, version 4.03,
V, Theorem 3.6, Corollary 3.7 and Example 3.9, pp.~158--159,
\url{https://www.jmilne.org/math/CourseNotes/CFT.pdf}.
The norm compatibility can also be checked directly on prime ideals
using the residue degree in the Frobenius formula.

\item[TW-CD: abelian conductor--discriminant formula.]
For an abelian number field, its absolute discriminant is the product
of the conductors of its characters.  In particular,
$F\subseteq\mathbb Q(\zeta_N)$ implies
$|\Delta_F|\le N^{[F:\mathbb Q]}$.
See Sharifi, equation (4.2.1),
\url{https://math.ucla.edu/~sharifi/notes/iwasawa-ch04.html}.
This is the same classical discriminant input already recorded in
OB-MD.
\end{description}

Finite Poitou--Tate duality is the already reconstructed LD3, rather
than another unnamed classical assumption.  LK2 and LK3 realize its
finite local-cohomology and local-duality ingredients on the actual
finite extensions inside the specified $\ell$-adic algebraic closure.
Elementary finite abelian groups, compact duality, the two-term
cohomology complex of a procyclic pro-$\ell$ group, and the elementary
ramification structure of local fields are used explicitly below.

\paragraph{TW1: the exact Ferrero--Washington rank interface, including two.}
\label{claim:TW1}
Let $F/\mathbb Q$ be abelian and contain $\mu_\ell$ if $\ell$ is odd,
or $\mu_4$ if $\ell=2$.  There exists an integer $r\ge0$ such that
for every $a\ge0$ and every field model $F'$ of
$F(\mu_{\ell^a})$ over $F$,
\[
 \left|\operatorname{Cl}(F')/\ell\operatorname{Cl}(F')\right|
       \le\ell^r.
 \tag{TW.3}
\]
Moreover, the continuous everywhere-unramified characters
$G_{F(\mu_{\ell^\infty})}\to\mathbb F_\ell$ have dimension at most
one fixed finite integer.  Both statements are invariant under
changing the field model or the algebraic closure.

\begin{proof}
The cyclotomic image of $G_F$ is open in
$\mathbb Z_\ell^\times$.  The root hypothesis puts it inside
$1+\ell\mathbb Z_\ell$ for odd $\ell$, and inside
$1+4\mathbb Z_2$ at two.  The $\ell$-adic logarithm, or the binomial
valuation formula, identifies these groups with $\mathbb Z_\ell$.
Their open subgroups are $1+\ell^c\mathbb Z_\ell$, with $c\ge1$
and $c\ge2$ at two.  Consequently
$F_\infty=F(\mu_{\ell^\infty})$ is precisely the cyclotomic
$\mathbb Z_\ell$-extension of $F$; its root-adjoining fields are
the usual finite layers, possibly repeated at the beginning.
In particular, no extra order-two quotient is present at $\ell=2$.

By TW-IW and TW-FW, the inverse-limit class module $X$ is finitely
generated over $\mathbb Z_\ell$ itself.  Indeed its structure theorem
has no $\Lambda/(\ell^b)$ terms, and every remaining
$\Lambda/(f^a)$, with $f$ distinguished, is finite free over
$\mathbb Z_\ell$.  Finite kernel and cokernel preserve finite
generation over the noetherian ring $\mathbb Z_\ell$.

We spell out why this gives the actual finite-layer assertion
(TW.3), without assuming a stronger control theorem.
At a prime above $\ell$ the local cyclotomic extension is eventually
totally ramified.  To check this after base change, let $K/\mathbb
Q_\ell$ be the local base and let $L_b/\mathbb Q_\ell$ be a finite
cyclotomic root layer.  It is totally ramified, while
$[KL_b:K]\le[L_b:\mathbb Q_\ell]$ and
$e(KL_b/K)\ge e(L_b/\mathbb Q_\ell)/e(K/\mathbb Q_\ell)$.
Consequently the residue degrees $f(KL_b/K)$ are bounded by
$e(K/\mathbb Q_\ell)$; in the increasing tower they stabilize.
The local cyclotomic image is open in the global
$\mathbb Z_\ell$-image, so the number of primes also stabilizes.
Thus, after one layer $n_0$,
every subsequent extension of layers is totally ramified at a
chosen prime above $\ell$.  If $m\ge n\ge n_0$, no nontrivial
subextension of $F_m/F_n$ is unramified at that prime.
Let $H_n/F_n$ be the Hilbert $\ell$-class field.  Then
$H_n\cap F_m=F_n$, and $H_nF_m/F_m$ is everywhere unramified.
It is an abelian $\ell$-extension, so it lies in $H_m$.
Restriction therefore maps
$\operatorname{Gal}(H_m/F_m)$ onto
$\operatorname{Gal}(H_n/F_n)$.  By TW-CF this is the norm
$\operatorname{Cl}(F_m)\{\ell\}\to
\operatorname{Cl}(F_n)\{\ell\}$.

All these later norm maps are surjective.  Compactness of the
inverse limit of finite groups gives a surjection
$X\to\operatorname{Cl}(F_n)\{\ell\}$ for every $n\ge n_0$.
Hence their mod-$\ell$ dimensions are at most
$\dim_{\mathbb F_\ell}X/\ell X$.  Take the maximum with the
finitely many earlier dimensions.  The prime-to-$\ell$ part of a
finite class group has zero quotient by $\ell$, proving (TW.3).
This is the literal statement
\texttt{Published.FerreroWashington}.

The class-field identification in TW-IW gives
\[
 \operatorname{Hom}_{\mathrm{cont,unr}}
 (G_{F_\infty},\mathbb F_\ell)
   =\operatorname{Hom}_{\mathrm{cont}}(X,\mathbb F_\ell),
 \tag{TW.4}
\]
so the same finite generator bound suffices at infinity.
Finally, a field isomorphism transports ideals, class groups, roots
of unity and their adjoining fields.  Extending it to algebraic
closures conjugates the corresponding absolute Galois groups and
their inertia subgroups.  All maps just used therefore preserve
both dimensions, for every field and closure model.
\end{proof}

\paragraph{TW2: one two-primary rank bound before the auxiliary primes are chosen.}
\label{claim:TW2}
Fix the abelian field datum $M$ with modulus $m>0$ and degree
$o=M.\deg$.  There is an integer $R_0(M)\ge1$, independent of
$q,d$, with the following property.  For every \texttt{PrimePair}
$(q,d)$ there is an abelian field $F_0$ such that
\[
 \mu_4\subset F_0,\quad [F_0:\mathbb Q]\le8o,\quad
 F_0\hookrightarrow\mathbb Q(\zeta_{4mqd}),
 \tag{TW.5}
\]
and the compact Galois module of the maximal unramified abelian
pro-$2$ extension of $F_0(\mu_{2^\infty})$ in which primes above
$2$ split is generated by at most $R_0(M)$ elements over
$\mathbb Z_2$.  One may take
\[
                 R_0(M)=4(g+6o),                         \tag{TW.6}
\]
where $g$ is a fixed bound from TW1 for $M(i)$.
For every $s\ge2$ and every \emph{original} field $E$ satisfying
$\texttt{IsCompField}\ M\ s\ q\ d\ E$, some layer $F_a/F_0$ with
$a\le s+2$ satisfies, for every $k\ge1$,
\[
                    s_2(E,k)\le1+s_2(F_a,k).             \tag{TW.7}
\]
For each odd prime $\ell$ and abelian datum $A$ containing
$\mu_\ell$, an analogous fixed finite generator bound
$R_\ell(A)\ge1$ exists for every field model of $A$.

\begin{proof}
Choose a model $M_0$ of $M$ inside $\overline{\mathbb Q}$.
For an odd prime $p$, let $Q_p$ be the quadratic subfield of
$\mathbb Q(\zeta_p)$ fixed by the squares in
$(\mathbb Z/p)^\times$.  It is unramified outside $p$.
Set
\[
 H_\infty=M_0(\mu_{2^\infty}),\qquad
 F_0=M_0(i)Q_qQ_d,\qquad
 F_\infty=H_\infty Q_qQ_d.                              \tag{TW.8}
\]
These fields give (TW.5) immediately.  The root and degree data
are uniform even if some composita have smaller than their maximal
possible degree.

Let $\mathcal U$ be the $\mathbb F_2$-space of characters of
$G_{H_\infty}$ unramified outside primes above $q,d$.
Its everywhere-unramified subspace has dimension at most $g$
by TW1.  Restriction to the inertia groups at those primes has
that subspace as its kernel.  At an odd residue characteristic,
wild inertia has no nontrivial map to $\mathbb F_2$, and the
maximal pro-$2$ quotient of tame inertia is procyclic.
Thus each prime contributes at most one dimension.

There are at most $2o$ primes of $H_\infty$ over $q$ and at most
$4o$ over $d$.  Here is the precise congruence calculation.
For $q\equiv3\pmod8$, one has $v_2(q^2-1)=3$, so for $n\ge3$
the order of $q$ in $(\mathbb Z/2^n)^\times$ is $2^{n-2}$.
For $d\equiv7\pmod{16}$, one has $v_2(d^2-1)=4$, and for
$n\ge4$ its order is $2^{n-3}$.  These statements follow by
successively squaring $1+2^b u$ with $u$ odd; each squaring
increases its distance valuation from $1$ by exactly one.
The Frobenius indices in the respective cyclotomic fields are
therefore $2$ and $4$.  Passing from $\mathbb Q(\mu_{2^\infty})$
to its compositum with $M_0$ multiplies the number of primes by
at most $o$.  The bounded increasing finite-layer prime counts
stabilize, so these are also the counts at infinity.  We obtain
\[
                     \dim\mathcal U\le g+6o.            \tag{TW.9}
\]

For completeness, an index-two extension multiplies this bound
by at most two.  If $V\triangleleft U$ has quotient $C_2$, set
$W=\operatorname{Hom}_{\rm cont}(V,\mathbb F_2)$.
For the induced involution $\tau$, $(\tau-1)^2=0$, so
$\dim W\le2\dim W^{C_2}$: its image lies in its kernel.
The restriction--inflation sequence
\[
 0\longrightarrow H^1(C_2,\mathbb F_2)
 \longrightarrow H^1(U,\mathbb F_2)
 \longrightarrow W^{C_2}
 \longrightarrow H^2(C_2,\mathbb F_2)
 \tag{TW.10}
\]
and the fact that both outer groups have dimension one imply
$\dim W^{C_2}\le\dim H^1(U,\mathbb F_2)$.
This also proves finiteness before using the dimension formula.
Apply this twice to the extensions adjoining $Q_q$ and $Q_d$.
They are unramified outside the allowed primes, so their maximal
extensions unramified outside those primes can be taken inside
one common Galois closure; the indicated subgroups really have
index at most two.  A trivial extension only improves the bound.
The everywhere-unramified, $2$-split characters of $G_{F_\infty}$
form a subspace, so their dimension is at most (TW.6).

Let $X$ be the compact $\mathbb Z_2$-dual of the discrete group
of all such $\mathbb Q_2/\mathbb Z_2$-valued characters.
Then $(X/2X)^\vee$ is the space just bounded.  Compact Nakayama
gives a surjection $\mathbb Z_2^{R_0}\to X$.
One direct proof is to lift a basis of $X/2X$, approximate an
arbitrary element successively modulo $2,2^2,\ldots$, and take
the limit of its coefficients in $\mathbb Z_2$.
The image is compact and closed, and $\bigcap_n2^nX=0$:
an element of this intersection annihilates every finite-order
character in the discrete dual.  The subgroups $2^nX$ tend to
zero for the same reason.
The same argument applies to every prime $\ell$.

The fixing subgroup in \texttt{IsCompField} identifies the actual
image of $E$ with
$M_0\,\mathbb Q(\zeta_{2^{s+2}})^+ Q_qQ_d$ in its given
cyclotomic overfield.  Adjoining $i$ gives
$F_a=F_0(\mu_{2^{s+2}})$, whose degree over $F_0$ is $2^a$
for some $a\le s+2$.  The extension $F_a/E$ has degree at most
two.  Restriction of cohomology preserves every unramified and
strict local condition in (TW.1).  Its kernel embeds in
$H^1(C_2,A_2(k)^{G_{F_a}})$ when the degree is two.
The invariant coefficient group is either a finite cyclic
$2$-group or $\mathbb Q_2/\mathbb Z_2$.
For a cyclic coefficient subgroup, the formula
$H^1(C_2,B)=\ker(1+\tau)/(\tau-1)B$ shows that it is killed by
$2$ and has at most one generator; its order is at most $2$.
For the divisible group the only involutions are $1,-1$:
the same formula gives respectively $B[2]$ and $B/2B=0$.
Thus the kernel length is at most one, proving (TW.7).
Conjugating the embedding back transports this inequality to
the original $E$, not to a newly substituted host.

Finally, for a fixed odd $\ell$ and $A$, the condition
$A.\texttt{HasMu}\ \ell$ means its fixing subgroup acts trivially
on $\mu_\ell$, so every model contains those roots.
TW1 bounds the mod-$\ell$
dimension of its unramified tower module.  Taking the
$\ell$-split quotient cannot increase the number of generators.
Enlarge the bound to at least one, and use the field-model
transport at the end of TW1.  This supplies $R_\ell(A)$.
\end{proof}

\paragraph{TW3: exact twisted tower descent and the reference-layer comparison.}
\label{claim:TW3}
Let $F_0$ contain $\mu_\ell$ if $\ell$ is odd and $\mu_4$ at two.
Put $F_\infty=F_0(\mu_{\ell^\infty})$, choose a topological
generator $\gamma$ of $\Gamma=\operatorname{Gal}(F_\infty/F_0)$,
and put $u=\chi_{\rm cyc}(\gamma)$.
Let $X$ be the compact module of unramified abelian pro-$\ell$
extensions of $F_\infty$ split at primes above $\ell$,
with its conjugation action $\rho(\gamma)$.
For every $k\ge1$ and every model of the $a$th layer,
\[
 s_\ell(F_a,k)=
 \operatorname{length}_{\mathbb Z_\ell}
 \frac{X}{((u^k\rho(\gamma))^{\ell^a}-1)X}.              \tag{TW.11}
\]
For every original $\Omega,D$ on any such layer $F$, LD3 and
LK2--LK3 give
\[
                         s_\ell(F,k)\le V_\ell(F,k;D).
 \tag{TW.12}
\]
Both statements keep the negative twist in (TW.1).

\begin{proof}
Over $F_\infty$ the twist is trivial.  Restrict a global cocycle
to $G_{F_\infty}$; its restriction is a continuous additive
character.  The strict local conditions say precisely that it
vanishes on every inertia group and on each decomposition group
above $\ell$.  Let $D_\infty=X^\vee$ be this character group.
Each finite-image character has an open stabilizer under conjugation:
continuity on the compact Galois group gives one open neighborhood
working uniformly on its finitely many fibers.  Thus this is a
discrete continuous $\Gamma$-module, and its compact dual action
is continuous.  Conjugation on it is
$f(h)\mapsto f(\gamma^{-1}h\gamma)$; the twisted action is
\[
         f(h)\longmapsto u^{-k}f(\gamma^{-1}h\gamma).
 \tag{TW.13}
\]
Restriction therefore lands in the invariants for (TW.13)
under $\gamma^{\ell^a}$.

This restriction is an isomorphism.  Indeed the continuous
cohomology complex of $\mathbb Z_\ell$ on a discrete
$\ell$-primary module is the two-term complex
$A\xrightarrow{\gamma-1}A$.  In the present module
$u^{-k\ell^a}-1\ne0$, and multiplication by this scalar on
$\mathbb Q_\ell/\mathbb Z_\ell$ is surjective.
Thus $H^1(\Gamma^{\ell^a},A_\ell(k))=0$ and
$H^2(\Gamma^{\ell^a},A_\ell(k))=0$.
The restriction--inflation sequence proves injectivity and
surjectivity onto invariant characters.
The local conditions survive the inverse construction.
Away from $\ell$ the cyclotomic tower is unramified, so the
whole inertia subgroup lies in $G_{F_\infty}$ and the lifted
cocycle vanishes there.  At a prime above $\ell$, the local
cyclotomic image is open in $\Gamma^{\ell^a}$.
A class whose restriction to the corresponding infinite-level
decomposition kernel is zero is inflated from this open
procyclic group; its $H^1$ vanishes by the same nonzero-scalar
argument.  Hence the lift is locally zero as required.

The dual of (TW.13) has positive twist
$u^k\rho(\gamma)$.  Duality changes invariants into
coinvariants.  If one first obtains the inverse operator,
its image minus the identity is the same image, since
$g^{-1}-1=-g^{-1}(g-1)$ for an automorphism $g$.
Compact/discrete duality preserves length, including the value
$\infty$, giving (TW.11).  There is a unique subgroup of index
$\ell^a$ in $\Gamma$; thus every field satisfying
\texttt{IsLayer} is carried to this same layer by an embedding
extending the specified $F_0$-algebra structure.

We give the limit argument for (TW.12), rather than treating
it as an extra duality hypothesis.
For $r\ge1$ let
$T_r=\mathbb Z/\ell^r(k+1)$ and
$A_r=A_\ell(k)[\ell^r]$.
The Cartier dual of $T_r$ is exactly $A_r$:
the roots-of-unity twist subtracts $k+1$ from $1$, giving
$-k$.  The finite Poitou--Tate pairing in LD3 identifies
the duals of the finite groups
\[
 \operatorname{Sha}^2(F,T_r)=
 \ker\!\left(H^2(G_{F,S_\ell},T_r)\to
          \prod_{v\mid\ell}H^2(F_v,T_r)\right),
 \qquad \operatorname{Sha}^1(F,A_r).
 \tag{TW.14}
\]
These are the natural cup-product pairings.  Reduction
$T_{r+1}\to T_r$ is paired with inclusion
$A_r\to A_{r+1}$: this follows directly by evaluating the
coefficient pairing and then the same local invariant maps.
The local transport proved in KB6 identifies the original
$\operatorname{Dec}(D.A_v)$ with the absolute Galois group of
the corresponding henselization, then of its completion $F_v$,
and finally of a finite field
$K_v\subset\texttt{PadicAlgCl}(\ell)$.
These are continuous isomorphisms preserving every root of
unity and hence the original cyclotomic character.
Pullback on cochains preserves cup products and the given
localization map.  LK2--LK3 therefore supply those maps and
finite local cohomology on the prescribed local fields,
including at $\ell=2$ and at every coefficient level $r$.
There are no real-place terms, since the root hypothesis
implies that $F$ is totally complex.

All finite-level $H^1(G_{F,S_\ell},T_r)$ are finite, as also
established on these exact cochains in KB1:
their localization kernels are finite by finite duality,
applied with the Cartier-dual coefficient as necessary,
and their images lie in the finite product of the finite
local $H^1$ groups.  Their inverse systems therefore satisfy
the Mittag--Leffler condition.  Continuous cochains with
coefficients in $\mathbb Z_\ell(k+1)$ are the inverse limits
of the finite coefficient cochains; transition maps on cochains
are surjective by lifting a function on its finitely many
clopen fibers.  The usual kernel--cokernel calculation for
inverse-limit complexes, with the just established
$\varprojlim{}^1H^1=0$, gives
$H^2(F,T_\ell(k))=\varprojlim_r H^2(F,T_r)$, globally
and locally.  Inverse limits preserve kernels, so
$\operatorname{Sha}^2(F,T_\ell(k))
=\varprojlim_r\operatorname{Sha}^2(F,T_r)$.

On the torsion side continuous cohomology commutes with this
filtered union of coefficients: every continuous cochain on a
compact profinite group has finite image.  There are only
finitely many local terms, and filtered unions are exact.
Consequently
$\operatorname{Sha}^1(F,A_\ell(k))
=\varinjlim_r\operatorname{Sha}^1(F,A_r)$.
Taking the limit of the compatible finite pairings in (TW.14)
identifies its compact dual with $\operatorname{Sha}^2(F,T_\ell(k))$.
This latter group is a submodule of the actual global $H^2$
in (TW.2), which proves (TW.12).  All these constructions
commute with field isomorphisms and with conjugating a chosen
place in $\Omega$; therefore (TW.12) holds for every original
local-data record, not just a convenient representative.
\end{proof}

\paragraph{TW4: complete coinvariant bounds with at most \(R\) centers.}
\label{claim:TW4}
Suppose the module $X$ in TW3 is generated by $R\ge1$ elements
over $\mathbb Z_\ell$.  Set
\[
 j_0=\lfloor\log_\ell R\rfloor+3,\qquad
 c=v_\ell(u-1).
 \tag{TW.15}
\]
There are $s_1,\ldots,s_R\in\mathbb Z_\ell$ such that
for every $k\ge1$ and every $a\ge0$,
\[
 s_\ell(F_a,k)\le
   \sum_{i=1}^R\bigl(\max(a,j_0)+c+v_\ell(s_i-k)\bigr).
 \tag{TW.16}
\]
Here $v_\ell(0)=\infty$, and
$c\ge1$, with $c\ge2$ at two; moreover $F_0$ contains a
primitive $\ell^c$th root of unity.
For every model $F_{j_0}$ of the reference layer and every
local-data record $D$ on it,
\[
 s_\ell(F_a,k)\le
 V_\ell(F_{j_0},k;D)+R\max(a-j_0,0).                    \tag{TW.17}
\]
These are the two parts of the literal \texttt{ABound l R F0}.

\begin{proof}
The cyclotomic image is $1+\ell^c\mathbb Z_\ell$ by the
root calculation in TW1, and $c$ is the largest exponent for
which $\mu_{\ell^c}\subset F_0$.
The action of $\Gamma$ on the finite vector space $X/\ell X$
is continuous, hence factors through a finite $\ell$-group.
Thus $\rho(\gamma)-1$ is nilpotent modulo $\ell$, with
nilpotence degree at most $\dim X/\ell X\le R$.
For $g=u^k\rho(\gamma)$ the reduction is the same,
so this assertion is uniform in $k$.

We first prove the layer growth bound
\[
 \operatorname{length}\operatorname{coker}(g^{\ell^a}-1)
 \le
 \operatorname{length}\operatorname{coker}(g^{\ell^{j_0}}-1)
                +R\max(a-j_0,0).                       \tag{TW.18}
\]
Put $h=\lfloor\log_\ell R\rfloor+1$, so $\ell^h\ge R$.
In characteristic $\ell$,
$(g-1)^{\ell^h}=g^{\ell^h}-1$, hence
$(g^{\ell^h}-1)X\subset\ell X$.
One further $\ell$th power gives image in $\ell^2X$,
by the binomial theorem; this is valid at two as well.
It remains true after further $\ell$th powers.
If $(w-1)X\subset\ell^2X$, set
$\Phi_w=1+w+\cdots+w^{\ell-1}$.
For every $x$, $\Phi_wx=\ell x+\ell^2y$ for some $y$.
Thus $\ell X\subset\Phi_wX+\ell(\ell X)$.
Nakayama applied to
$(\ell X+\Phi_wX)/\Phi_wX$ gives
$\ell X\subset\Phi_wX$, so
$\operatorname{length}\operatorname{coker}\Phi_w\le R$.
Factor $w^\ell-1=\Phi_w(w-1)$ and use the exact sequence
of cokernels of a composite; the length increases by at most
$R$ at each later layer.  At earlier layers monotonicity follows
from the same factorization of $g^{\ell^{j_0}}-1$ by
$g^{\ell^a}-1$.  This proves (TW.18), also when a length is
infinite.

Choose a surjection $\pi:\mathbb Z_\ell^R\to X$, and lift
$\rho(\gamma)^{\ell^{j_0}}$ to a matrix $S$ on its generators.
Then $\pi S=\rho(\gamma)^{\ell^{j_0}}\pi$.
There is a surjection of the corresponding cokernels, so
Smith normal form gives
\[
 \operatorname{length}\operatorname{coker}
       (g^{\ell^{j_0}}-1)
 \le v_\ell\det\bigl((u^{\ell^{j_0}})^kS-I\bigr).
 \tag{TW.19}
\]
A zero determinant gives $\infty$, consistently with the
inequality.  In an algebraic extension of $\mathbb Q_\ell$
factor the determinant as
$\prod_{i=1}^R(a_k z_i-1)$, where
$a_k=(u^{\ell^{j_0}})^k$ and the $z_i$ are the eigenvalues
of $S$, with multiplicities.

For each $z=z_i$, we construct a center in $\mathbb Z_\ell$
without requiring $z$ to be a power of $u$.
For all nonnegative integers $b,k$ the binomial valuation
formula gives
\[
 |a_b-a_k|_\ell=\ell^{-j_0-c}|b-k|_\ell.
 \tag{TW.20}
\]
For any units $A,B$ and any $z$ in the normed extension,
\[
 |A-B|_\ell\le
       \max\{|Az-1|_\ell,\ |Bz-1|_\ell\}.
 \tag{TW.21}
\]
If $|z|_\ell\ge1$, this follows by subtracting the two
expressions and dividing by $|z|_\ell$.
If $|z|_\ell<1$, both norms on the right are one.
Therefore the closed balls in $\mathbb Z_\ell$ centered at
$k$ with radii $\ell^{j_0+c}|a_kz-1|_\ell$ intersect
pairwise.  Intersecting ultrametric balls are nested; their
finite intersections are nonempty.  Compactness of
$\mathbb Z_\ell$ gives a common point $s_i$.
It satisfies
$v_\ell(a_kz_i-1)\le j_0+c+v_\ell(s_i-k)$.
Summing these inequalities in (TW.19), then using (TW.18)
and TW3, proves (TW.16).

Finally, (TW.18) and (TW.11) compare the arbitrary layer with
the reference layer.  Apply (TW.12) to that reference layer
with its given $D$, giving (TW.17).
The centers may be regarded as elements of $\mathbb Q_\ell$,
which is an allowed normed extension in \texttt{ABound}.
Taking $R'=R$ in its first part and translating
$v_\ell(s_i-k)\le B$ to
$\ell^{-B}\le|s_i-k|_\ell$ matches its norm convention exactly.
\end{proof}

\paragraph{TW5: exceptional twists and a common square window.}
\label{claim:TW5}
Under TW4, for every integer $Z\ge1$ there is a set
$\mathcal E\subseteq\{1,\ldots,Z\}$ of size at most $R$
such that, simultaneously for every layer and every
$1\le k\le Z$ outside $\mathcal E$,
\[
 s_\ell(F_a,k)\le
 R\bigl(\max(a,j_0)+c+\lfloor\log_\ell Z\rfloor\bigr).
 \tag{TW.22}
\]
Independently, let $\mathcal E$ be any finite set of
nonnegative integers of size at most $R$, let
$0<\eta\le C_t$, and let $X>1$.  With
\[
                 J=2^{R+2}(C_t/\eta)^R,                \tag{TW.23}
\]
there is $Y$ such that
\[
 X^{1/J}\le Y,\qquad Y^2\le X,\qquad
 \left.
 \begin{array}{c}
 x\in[Y,Y^2],\ k\in\mathcal E,\ k\le x^{C_t}
 \end{array}\right\}
       \Longrightarrow k\le x^\eta .                  \tag{TW.24}
\]

\begin{proof}
For each center $s_i$ select the integers $1\le k\le Z$
with $|s_i-k|_\ell<\ell^{-\lfloor\log_\ell Z\rfloor}$.
There is at most one: two distinct choices would have
$\ell$-adic difference valuation at least
$\lfloor\log_\ell Z\rfloor+1$, so their ordinary integer
difference would be at least
$\ell^{\lfloor\log_\ell Z\rfloor+1}>Z$.
Take the union of these sets and apply (TW.16).
This proves the full \texttt{FirstBound}, not merely a count
of exceptional twists.

For the second assertion put $B=C_t/\eta\ge1$ and
$U=\log X>0$.  We use the following elementary interval fact:
if a finite real set $T$ has at most $R$ elements, some $y$
satisfies
\[
 \frac{U}{2^{R+2}B^R}\le y,\quad 2y\le U,\quad
 a\in T,\ y\le v\le2y,\ a\le v\ \Longrightarrow\ Ba\le v.
 \tag{TW.25}
\]
Here is an induction proof, including its constants.
For $R=0$, take $y=U/2$.
If there is an $a_0\in T$ with $U/(2B)<a_0\le U$,
remove it and apply the induction hypothesis at
$U'=U/(2B)$ and size bound $R-1$.
Its lower bound is exactly the one in (TW.25), while
$v\le2y\le U'<a_0$ makes the deleted point irrelevant.
If no such point exists, take $y=U/2$.
Every $a\in T$ with $a\le v\le U$ then satisfies
$a\le U/(2B)$, so $Ba\le U/2\le v$.
The required lower bound is weaker than $U/2$.

Apply (TW.25) to
$T=\{\log k/C_t:k\in\mathcal E,\ k\ge1\}$ and take
$Y=\exp y$.  If $k\le x^{C_t}$, then
$a=\log k/C_t\le\log x$.  The conclusion
$Ba\le\log x$ is exactly $k\le x^\eta$.
The value $k=0$ causes no difficulty.  Exponentiating the
other inequalities proves (TW.24).  A union of exception
sets can be used here once, yielding the same $Y$ for all
of them.
\end{proof}

\paragraph{TW6: uniform reference-field sizes and exact model transfer.}
\label{claim:TW6}
Fix $\ell,j,n_1,N_0$ with $\ell$ prime and $N_0\ge1$.
There exist $n_0\ge1$ and $c_0\ge0$, depending only on these
four parameters, such that the following holds.
If $[F_0:\mathbb Q]\le n_1$,
$F_0\hookrightarrow\mathbb Q(\zeta_{N_0P})$, $P\ge1$,
$F_j/F_0$ is a $j$th layer, and
$1\le w$, $\log P\le w$, then
\[
 [F_j:\mathbb Q]\le n_0,\qquad
             \log(2|\Delta_{F_j}|)\le c_0w.             \tag{TW.26}
\]
For a fixed abelian datum $A$ and a fixed layer degree
$\ell^j$, all models of $A$ and all such layers are isomorphic
as towers.  The lengths (TW.1) are preserved by these
isomorphisms for every closure and local-data record.

\begin{proof}
Take $n_0=(n_1+1)\ell^j$.
Write $F_j=F_0(\mu_{\ell^b})$ for some $b$.
Since $F_j$ contains a primitive $\ell^b$th root,
$\varphi(\ell^b)\le[F_j:\mathbb Q]$.
For a prime power this implies
$\ell^b\le2[F_j:\mathbb Q]\le2n_0$, also for $b=0$.
Embedding the compositum in
$\mathbb Q(\zeta_{\ell^bN_0P})$ and using TW-CD gives
\[
 \log(2|\Delta_{F_j}|)
 \le \log2+n_0\log(2n_0N_0)+n_0\log P.
 \tag{TW.27}
\]
Thus one may take
$c_0=\log2+n_0\log(2n_0N_0)+n_0$.
The degree bound is the tower formula
$[F_j:\mathbb Q]=[F_0:\mathbb Q]\ell^j$.

A model of $A$ is the fixed field of its specified subgroup
of a cyclotomic Galois group, so two models have a rational
isomorphism.  Extend it to algebraic closures.  It transports
roots of unity and hence the root-adjoining towers.  These
fields are nested, so two with the same relative degree are
equal; under the hypothesis of TW1 this is also the unique
index-$\ell^j$ subgroup description.
The resulting isomorphism respects the
given base-algebra maps.  Conjugation transports absolute
Galois groups, the cyclotomic character, every inertia group
and every decomposition group.  Pullback on continuous
cocycles and its inverse commute with coboundaries and
localization.  Thus they give a linear isomorphism of the
strict kernels (TW.1).  Changing which place above a finite
prime represents a local term only conjugates that
restriction map, inducing the same kernel on cohomology.
This proves the stated model and local-data invariance.
\end{proof}

\paragraph{TW7: all twists on a good interval, from the literal \(\texttt{LemmaK}\).}
\label{claim:TW7}
Suppose $C_K\ge1$ satisfies the literal \texttt{LemmaK}:
for every prime $\ell$, every totally complex number field $F$
containing the roots required in TW1, every original $\Omega,D$,
and every $k\ge1$, the actual compact-coefficient group is
finite and
\[
 \log\left|H^2(G_{F,S_\ell},\mathbb Z_\ell(k+1))\right|
 \le
 \exp\!\left((2k[F:\mathbb Q])^{C_K}
                         \log(2|\Delta_F|)\right).       \tag{TW.28}
\]
Fix the prime $\ell$ and $R,c_1,n_0,c_0,J,C_t,\varepsilon,a$ with
$R\ge1$, $c_0,a,C_t\ge0$, $J,\varepsilon>0$.
There is $C\ge0$ depending only on these fixed parameters
and $C_K$ with the following property.
Suppose a tower as in TW4 has $c\le c_1$, its exception set
from TW5 is constructed up to $Z=\lfloor X^{C_t}\rfloor$,
and its reference layer satisfies
\[
 [F_{j_0}:\mathbb Q]\le n_0,\qquad
 \log(2|\Delta_{F_{j_0}}|)\le c_0x^{\varepsilon/4}.        \tag{TW.29}
\]
If $1\le x\le X$, $X^{1/J}\le x$, and
\[
 k\in\mathcal E,\ k\le x^{C_t}
              \Longrightarrow k\le x^{\varepsilon/(2C_K)},
 \tag{TW.30}
\]
then for every $1\le k\le x^{C_t}$ and every layer index
$b\le a(1+\log x)$ the length $s_\ell(F_b,k)$ is finite and
\[
 s_\ell(F_b,k)+1\le
       C(1+\log x)+2\exp(Cx^{3\varepsilon/4}).
 \tag{TW.31}
\]
For fixed $C,C_t,\varepsilon>0$, the right side times $x^{C_t}$
is at most $\exp(x^{4\varepsilon/5})$ for all sufficiently
large $x$.

\begin{proof}
A finite $\mathbb Z_\ell$-module of length $v$ has order
$\ell^v$.  Since $\log\ell\ge\log2>1/2$, (TW.28) gives
\[
 V_\ell(F,k;D)\le
 2\exp\!\left((2k[F:\mathbb Q])^{C_K}
                                  \log(2|\Delta_F|)\right).
 \tag{TW.32}
\]
If $k\in\mathcal E$, (TW.30) and (TW.29) imply
\[
 (2k[F_{j_0}:\mathbb Q])^{C_K}
       \log(2|\Delta_{F_{j_0}}|)
 \le (2n_0)^{C_K}c_0x^{3\varepsilon/4}.                 \tag{TW.33}
\]
Apply (TW.17) and (TW.32).  If $k\notin\mathcal E$, use
(TW.22); the inequality $k\le Z$ follows from
$k\le x^{C_t}\le X^{C_t}$.
Both cases prove finiteness and the common bound
\[
 s_\ell(F_b,k)\le
 R\bigl(b+j_0+c+\lfloor\log_\ell Z\rfloor\bigr)
       +2\exp((2n_0)^{C_K}c_0x^{3\varepsilon/4}).         \tag{TW.34}
\]
The logarithmic inequalities
$\log X\le J\log x$ and
$\log\ell>1/2$ give
$\lfloor\log_\ell Z\rfloor\le2C_tJ\log x$.
Thus (TW.31) holds with
\[
 C=R(a+j_0+c_1+2C_tJ)+1+(2n_0)^{C_K}c_0.
 \tag{TW.35}
\]
In particular the constant is uniform in the varying base
field whenever the displayed degree, discriminant and root
bounds are uniform.

For the final absorption, divide the logarithm of
$x^{C_t}[C(1+\log x)+2\exp(Cx^{3\varepsilon/4})]$
by $x^{4\varepsilon/5}$.
The exponential term contributes at most a constant times
$x^{-\varepsilon/20}$; the remaining terms are bounded by
constants times $\log x$ and $\log(2+\log x)$ divided by
$x^{4\varepsilon/5}$.  They all tend to zero.
The ratio is therefore at most one after a threshold depending
only on $C,C_t,\varepsilon$.  This proves the assertion.
\end{proof}

\paragraph{TW8: the literal simultaneous \(\texttt{ThmBab}\), with its quantifiers.}
\label{claim:TW8}
Assume $\exists C_K,\ \texttt{LemmaK}\ C_K$.
Together with TW1--TW7, LD3 and LK2--LK3, this supplies
\[
                \exists K:\texttt{BConst},\
                \forall C_t\ge1,\ \texttt{ThmBab}\ K\ C_t .
 \tag{TW.36}
\]
More explicitly, the single $K=(R_0,R_\ell,C_K)$ is chosen
before $C_t$, the abelian datum $M$, the finite set $B'$ of
odd primes, the family of abelian data $L(\ell)$ containing
$\mu_\ell$, $\varepsilon,A_1$, the admitted functions, the
prime pair and the eventual window.  For these data put
\[
 \begin{split}
 R&=R_0(M)+\sum_{\ell\in B'}R_\ell(L(\ell)),\\
 J&=2^{R+2}\left(\frac{2C_KC_t}{\varepsilon}\right)^R.
 \end{split}                                           \tag{TW.37}
\]
For $0<\varepsilon<1/52$ and any real $A_1$ there exists $X_0$
such that for every pair of functions $e,t:\mathbb R\to\mathbb N$
satisfying, for all $x\ge2$,
\[
 2\le e(x),\quad 2^{e(x)}\le x^{A_1},\quad
                  2\le t(x)\le x^{C_t},                 \tag{TW.38}
\]
every $X\ge X_0$, and every \texttt{PrimePair} $(q,d)$ with
\[
                         \log(qd)\le X^{\varepsilon/(4J)},
 \tag{TW.39}
\]
one common $Y$ satisfies $X^{1/J}\le Y$, $Y^2\le X$ and the
following two assertions for every $x\in[Y,Y^2]$:
\[
 \begin{array}{ll}
 \operatorname{VV}_2(E,t(x),\exp(x^{4\varepsilon/5}))
 &\text{for every original }
 E\text{ with }\texttt{IsCompField}\ M\ e(x)\ q\ d\ E;\\[2mm]
 \operatorname{VV}_\ell(F,t(x),\exp(x^{4\varepsilon/5}))
 &\text{for every }\ell\in B',\ \text{every model }F_0
   \text{ of }L(\ell),\\
 &\text{and every }j_0(\ell,R_\ell(L(\ell)))\text{th layer }
 F/F_0 .
 \end{array}                                           \tag{TW.40}
\]
The second line quantifies separate reference-layer models;
it does not identify them with the field $E$ in the first.

\begin{proof}
First enlarge a witness for \texttt{LemmaK} to $C_K\ge1$.
This preserves (TW.28), since $2k[F:\mathbb Q]\ge2$ and
$\log(2|\Delta_F|)>0$.  Choose the function $R_0(M)$ by
TW2 for every abelian datum.  For each odd $\ell$ and
root-containing datum $A$, choose its fixed bound
$R_\ell(A)$ by TW2; assign arbitrary values, for example zero,
outside that domain.  These choices and $C_K$ define one
\texttt{BConst}.  They are made before any $C_t$.
Now fix the remaining parameters in the statement.

For the two-primary part, write $o=M.\deg$,
$R_2=R_0(M)$, and $j_2=j_0(2,R_2)$.
Every $F_0$ obtained in TW2 has degree at most $8o$ and is
contained in modulus $4mqd$.  It contains a primitive
$2^c$th root.  Since
$\varphi(2^c)\le[F_0:\mathbb Q]$, we have the coarse uniform
bound $c\le16o$.  Apply TW6 with
$\ell=2$, $j=j_2$, $n_1=8o$, $N_0=4m$, and $P=qd$.
It gives $n_0,c_0$ independent of $q,d$.
If $x\ge X^{1/J}$, (TW.39) implies
$\log(qd)\le x^{\varepsilon/4}$, so every reference layer
in this tower satisfies (TW.29).

By (TW.38),
$e(x)\log2\le A_1\log x$, hence
$e(x)+2\le(2|A_1|+2)(1+\log x)$.
For the layer used in (TW.7), the parameter $a$ in TW7 may
therefore be fixed as $2|A_1|+2$.
TW7, with $c_1=16o$ and the uniform $n_0,c_0$, gives one
constant $C_2$ and an absorption threshold $x_2$ for this
whole family, before the functions, $X$ and $(q,d)$ are
chosen.  It follows that, whenever the exception set of this
tower is good as in (TW.30), each relevant length on the
original $E$ is bounded by the right side of (TW.31):
the extra $1$ in that bound accounts exactly for (TW.7).
There are $t(x)-1\le x^{C_t}$ twists.  Their sum is bounded
by $\exp(x^{4\varepsilon/5})$ once $x\ge x_2$.

For each of the finitely many $\ell\in B'$, choose one model
of $L(\ell)$ and its $j_0(\ell,R_\ell(L(\ell)))$th layer.
Its degree, discriminant and maximal root exponent are now
fixed constants.  Apply TW7 with that root exponent as $c_1$,
the exact degree as $n_0$, and
$c_0=\log(2|\Delta|)$; inequality (TW.29) holds for every
$x\ge1$.  Take $a=j_0$, so the fixed layer index obeys
$j_0\le a(1+\log x)$.
This gives a constant $C_\ell$ and a threshold $x_\ell$,
again before $e,t,X,q,d$.  TW6 transports every conclusion
to every other model of this base and every prescribed layer
over it.  It also transports finiteness and the full family
of original local-data records, not just numerical degrees.

Choose $x_*\ge2$ larger than all the finitely many thresholds
$x_2,x_\ell$, and then set $X_0=\max(2,x_*^J)$.
These choices have precisely the permitted dependencies.
Now choose arbitrary admitted $e,t$, $X\ge X_0$ and a prime
pair satisfying (TW.39).  Use TW2 to select its two-primary
base.  For this tower and for each fixed odd reference tower,
use TW5 with the same cutoff $Z=\lfloor X^{C_t}\rfloor$.
Let their exception sets be $\mathcal E_2,\mathcal E_\ell$.
Their union has size at most the $R$ in (TW.37).
Use the square-window part of TW5 once, with
$\eta=\varepsilon/(2C_K)$ and $C_t/\eta=
2C_KC_t/\varepsilon$.
The hypothesis $\eta\le C_t$ follows from
$C_K,C_t\ge1$ and $\varepsilon<1/52$.
It supplies exactly the $J$ in (TW.37) and one common $Y$.

For every $x\in[Y,Y^2]$, one has
$x\ge X^{1/J}\ge x_*$ and $x\le X$, and every constituent
exception set satisfies (TW.30).  The two preceding uniform
estimates therefore apply to all integers $1\le k<t(x)$.
Summing them proves both lines of (TW.40), for all the
specified fields and all their local data, on this same
window and with this fixed pair $(q,d)$.
By the definitions of \texttt{BaAt}, \texttt{BbAt},
\texttt{IsBWindow} and \texttt{ThmBab}, this is (TW.36).
\end{proof}

\subsection{Scope of the completed dependency chain}

TW1 supplies the source's exact
\texttt{Published.FerreroWashington} interface.
TW2--TW4 supply \texttt{TwoIn} and \texttt{FixedIn}, including
the complete \texttt{IsTower} to \texttt{ABound} step.
TW5--TW8 supply their simultaneous quantitative window, with one
$K$ before all $C_t\ge1$ and with all field models quantified.
The only project-specific incoming quantitative statement is
the literal \texttt{LemmaK}; KB8 supplies it with $C_K=31$.
The existing VW3--VW4 assembly can therefore replace its
\texttt{ThmBab} input by TW8 once that supplier is joined.
The odd-minus supplier is already OB8 with constant $240$.
The separate \texttt{PointHV} supplier for the odd-plus comparison,
and the other arithmetic, lattice, host and representation inputs
outside this branch, remain exactly where the cumulative dependency
audit records them.  No full \texttt{LinkN}, lattice gain, RH or
GRH conclusion is asserted by this chapter alone.
